\documentclass[10pt]{article}
\providecommand{\FIGDIR}{fig}
\newcommand{\DIGESTDATE}{10 August 2026}
% ===================== PM Research Watch - shared preamble =====================
\usepackage[T1]{fontenc}
\usepackage[utf8]{inputenc}
\usepackage{amsmath,amssymb}
\usepackage{tgpagella}
\usepackage{tgheros}
\renewcommand{\sfdefault}{qhv}
\renewcommand{\ttdefault}{lmtt}
\usepackage{microtype}
\usepackage[english]{babel}

\usepackage[a4paper,top=17mm,bottom=17mm,left=17mm,right=17mm,
            headsep=6mm,footskip=10mm]{geometry}
\usepackage{graphicx}
\usepackage{xcolor}
\usepackage{colortbl}
\usepackage{tikz}
\usetikzlibrary{calc}
\usepackage[most]{tcolorbox}
\usepackage{titlesec}
\usepackage{fancyhdr}
\usepackage{booktabs}
\usepackage{array}
\usepackage{longtable}
\usepackage{enumitem}
\usepackage{multicol}
\usepackage{ragged2e}
\usepackage[hidelinks]{hyperref}

% ------------------------------------------------------------------ palette
\definecolor{Ink}    {HTML}{14262E}
\definecolor{Deep}   {HTML}{0F4C5C}
\definecolor{Teal}   {HTML}{1F7A8C}
\definecolor{Sky}    {HTML}{4FA3C4}
\definecolor{Sage}   {HTML}{6FA287}
\definecolor{Amber}  {HTML}{E3A008}
\definecolor{Clay}   {HTML}{B5643C}
\definecolor{Coral}  {HTML}{D1495B}
\definecolor{Violet} {HTML}{6A5B9E}
\definecolor{Slate}  {HTML}{4A6572}
\definecolor{Paper}  {HTML}{FBF9F5}
\definecolor{Mist}   {HTML}{EEEAE2}
\definecolor{Rule}   {HTML}{D9D3C9}

\pagecolor{Paper}
\color{Ink}

\hypersetup{colorlinks=true, linkcolor=Deep, urlcolor=Teal, citecolor=Teal,
            pdftitle={PM Research Watch}, pdfauthor={Automated daily digest}}

% ------------------------------------------------------------------ headings
\titleformat{\section}
  {\sffamily\bfseries\large\color{Deep}}{}{0pt}
  {}[{\vspace{1.2mm}\color{Amber}\titlerule[1.1pt]}]
\titlespacing*{\section}{0pt}{5.5mm}{2.8mm}

\titleformat{\subsection}
  {\sffamily\bfseries\normalsize\color{Teal}}{}{0pt}{}
\titlespacing*{\subsection}{0pt}{3.6mm}{1.4mm}

\setlist[itemize]{leftmargin=4.2mm, itemsep=1.1mm, topsep=1.3mm, parsep=0pt,
                  label=\textcolor{Teal}{\footnotesize$\blacktriangleright$}}

% ------------------------------------------------------------------ header/footer
\pagestyle{fancy}
\fancyhf{}
\renewcommand{\headrulewidth}{0pt}
\renewcommand{\footrulewidth}{0pt}
\fancyhead[L]{\sffamily\scriptsize\color{Slate}PM RESEARCH WATCH\;\textcolor{Rule}{\textbar}\;\DIGESTDATE}
\fancyhead[R]{\sffamily\scriptsize\color{Slate}Particulate matter monitoring \& health effects}
\fancyfoot[C]{\sffamily\scriptsize\color{Slate}\thepage}

% ------------------------------------------------------------------ components
% metric tile
\newcommand{\metric}[3]{%
  \begin{tcolorbox}[enhanced, colback=white, colframe=Rule, boxrule=0.5pt,
    arc=1.6mm, left=2.2mm, right=2.2mm, top=1.8mm, bottom=1.8mm,
    borderline west={1.7mm}{0pt}{#3}, nobeforeafter, valign=center,
    height=17mm, width=\linewidth]
    \RaggedRight
    {\sffamily\bfseries\fontsize{17}{17}\selectfont\color{#3}#1}\\[0.5mm]
    {\sffamily\fontsize{6.5}{7.4}\selectfont\color{Slate}#2}
  \end{tcolorbox}}

% pull-quote / key insight box
\newtcolorbox{keybox}[1][]{enhanced, colback=Mist, colframe=Mist,
  arc=1.6mm, left=3.2mm, right=3.2mm, top=2.4mm, bottom=2.4mm,
  borderline west={1.7mm}{0pt}{Amber}, #1}

% subtopic band
\newcommand{\band}[2]{%
  \vspace{3.2mm}\noindent
  \begin{tcolorbox}[enhanced, colback=#2, colframe=#2, arc=1.2mm,
    left=2.6mm, right=2.6mm, top=1.1mm, bottom=1.1mm, boxrule=0pt,
    width=\linewidth, nobeforeafter]
    {\sffamily\bfseries\small\color{white}#1}
  \end{tcolorbox}\vspace{1.6mm}}

% one paper entry: {short cite}{title}{journal}{doi}{body}
\newcommand{\paperentry}[5]{%
  \par\vspace{1.9mm}\noindent
  {\sffamily\bfseries\footnotesize\color{Ink}#1}\;%
  {\sffamily\fontsize{7}{8}\selectfont\color{Slate}#3}\\[0.4mm]
  {\footnotesize\itshape\color{Deep}#2}\\[0.6mm]
  {\footnotesize #4}\\[0.3mm]
  {\fontsize{6.6}{7.6}\selectfont\color{Slate}\texttt{doi:}\,\href{https://doi.org/#5}{\texttt{\detokenize{#5}}}}
  \par}

% caption style
\newcommand{\figcap}[1]{\par\vspace{1.0mm}
  {\fontsize{7.2}{8.6}\selectfont\color{Slate}\textsf{\textbf{Fig.}}\; #1\par}\vspace{1.5mm}}

\providecommand{\permil}{\textperthousand}

\setlength{\parindent}{0pt}
\setlength{\parskip}{0pt}
\renewcommand{\arraystretch}{1.18}


\begin{document}

% ==================================================================== MASTHEAD
\thispagestyle{empty}
\begin{tikzpicture}[remember picture, overlay]
  \fill[Deep] (current page.north west) rectangle
      ($(current page.north east)+(0,-67mm)$);
  \fill[Amber] ($(current page.north west)+(0,-67mm)$) rectangle
      ($(current page.north east)+(0,-68.4mm)$);
  \foreach \x/\y/\r in {22/12/0.9, 41/26/1.5, 63/9/0.7, 88/31/1.1, 112/17/1.9,
                        138/29/0.8, 159/13/1.3, 178/54/1.0, 196/21/1.6,
                        30/44/0.6, 74/57/1.2, 125/48/0.9, 168/6/0.8, 191/59/0.7,
                        52/18/1.0, 100/60/0.6, 148/20/1.1, 66/49/1.4, 205/45/0.8}
    \fill[Sky, opacity=0.30] ($(current page.north west)+(\x mm,-\y mm)$) circle (\r mm);
\end{tikzpicture}

\vspace*{4mm}
{\sffamily\fontsize{7.6}{9}\selectfont\color{Sky}%
 AUTOMATED DAILY LITERATURE DIGEST\;\textbullet\;PhD RESEARCH WATCH\par}
\vspace{2.2mm}
{\sffamily\bfseries\fontsize{27}{29}\selectfont\color{white}%
 Particulate Matter\par}
{\sffamily\fontsize{27}{29}\selectfont\color{Sky}%
 Monitoring \& Health Effects\par}
\vspace{5.0mm}
{\sffamily\fontsize{8.4}{10}\selectfont\color{white}%
 Issue 2026-08-10\;\;\textcolor{Amber}{\textbar}\;\;Entry window 10 August 2026\;\;%
 \textcolor{Amber}{\textbar}\;\;16 records in scope (3 pre-watch backfill)\par}
\vspace{18mm}

% -------------------------------------------------------------- metric strip
\noindent
\begin{minipage}[t]{0.190\linewidth}\metric{16}{RECORDS IN SCOPE\\(27 REJECTED)}{Deep}\end{minipage}\hfill
\begin{minipage}[t]{0.190\linewidth}\metric{10}{SENSING / EXPOSURE\\RECORDS (63\%)}{Teal}\end{minipage}\hfill
\begin{minipage}[t]{0.190\linewidth}\metric{8}{SUBTOPIC CLUSTERS\\REPRESENTED}{Sky}\end{minipage}\hfill
\begin{minipage}[t]{0.190\linewidth}\metric{5}{METADATA-ONLY\\RECORDS}{Amber}\end{minipage}\hfill
\begin{minipage}[t]{0.190\linewidth}\metric{0}{RATIO-SCALE EFFECT\\ESTIMATES}{Clay}\end{minipage}

\vspace{6mm}

% -------------------------------------------------------------- signal of day
\section*{\sffamily\bfseries\color{Deep}Signal of the day}
\vspace{-1.5mm}{\color{Amber}\rule{\linewidth}{1.1pt}}\vspace{2.5mm}

\begin{keybox}
{\sffamily\bfseries\small\color{Deep}Two independent mobile-platform records land on the
same day, and between them they define the measurement problem this watch keeps
circling: dense mobile sampling produces the spatial detail that fixed networks cannot,
and produces it in a form nobody has agreed how to calibrate.}\\[1.4mm]
\footnotesize
Babu et al.\ (\textit{Environ Monit Assess}) drove a size-resolved optical instrument over
$\sim$2.4\,km$^2$ of Madurai at a 6\,s cadence on a 50\,m grid and recovered a
\textbf{bimodal} distribution with peaks at \textbf{16.77--32.37\,$\mu$m} and
\textbf{0.275--0.53\,$\mu$m}, local PM$_{10}$ enhancement of
\textbf{188.91\,$\pm$\,149.07\,$\mu$g/m$^3$} over background, and Getis-Ord Gi* hotspots
locked to traffic corridors. Bobbia et al.\ (\textit{Data in Brief}, pre-watch backfill
surfaced by the Amass connector) publish the complementary artefact: \textbf{10 mobile
AtmoTrack units on Rouen buses at 10\,s, 18 fixed AirSensEUR units on traffic lights at
1\,min, and 6 reference stations at 15\,min} over 10.5 months, assembled explicitly to
support \emph{multi-hop} calibration --- mobile sensor against fixed sensor against
reference --- because most mobile units never see a reference site.
Read together the implication is uncomfortable and specific: Babu's standard deviation is
\textbf{79\% of the mean enhancement}, and on a moving platform that dispersion is
inseparable from uncalibrated instrument response, sampling-line losses and 6\,s response
time. The Rouen dataset exists because that inseparability is the field's open problem;
the Madurai paper reports hotspots without addressing it. Neither is wrong. The gap
between them is where the calibration-transfer literature of the past fortnight ---
Ginsburg's regional PurpleAir corrections, Zhivkov's distance-dependent uncertainty ---
now has to go: from static co-location to platforms that move.
\end{keybox}

\vspace{3mm}

\clearpage

% -------------------------------------------------------------- exec brief
\section*{\sffamily\bfseries\color{Deep}What changed today}
\vspace{-1.5mm}{\color{Amber}\rule{\linewidth}{1.1pt}}\vspace{2.5mm}

\footnotesize
\begin{itemize}
\item \textbf{The connector layer was widened, and it changed what the day looks like.}
This issue is the first run to sweep the full set of connected research surfaces ---
PubMed, Amass BiomedCore, Consensus, Scholar Gateway, Semantic Scholar, arXiv, medRxiv,
Crossref, Exa and Nimble --- alongside the committed harvester. The measurable result is
\textbf{not} more same-day records: PubMed's entry-date window is still the only leg that
reliably resolves a single day. What the wider sweep produced was \textbf{three
high-value records with pre-watch entry dates} (9 July, 29 June, 9 June) that no
date-windowed query could ever have surfaced, and \textbf{seven duplicates} that
\texttt{seen.json} caught on PMID or DOI. Both outcomes are worth having; neither
substitutes for the entry-date spine.

\item \textbf{Instrument-side particle metrics finally outnumber mass metrics.}
Six of sixteen records report absorption coefficient, number size distribution,
heavy-metal content or optical apportionment rather than a regulated mass concentration.
The 9 August run log flagged that \texttt{pmclass()} had no bin for any of these and
collapsed them all into ``PM (unspeciated)''. Four new classes were added in this issue.
The fallback bin did not disappear --- it still holds \textbf{4} records --- but those
four now genuinely have no stated particle metric, and \textbf{5} records that would have
been buried there are separable as optical properties or composition/speciation.

\item \textbf{Source apportionment moves from composition to optics.}
Acton-Bond et al.\ apply \emph{rolling} positive matrix factorisation to aerosol
\emph{optical} apportionment rather than to filter-based chemical speciation. If the
approach holds, it puts source attribution within reach of an aethalometer rather than a
laboratory, which is directly relevant to what a distributed network can be asked to
report. The record is metadata-only, so this is a flag, not a finding.

\item \textbf{The measurement-comparability question is being asked out loud.}
Wang et al.\ title a paper ``operational definitions, method comparability and
environmental interpretation'' for heavy metals in PM. That is the same critique this
watch has been making of low-cost mass measurement, applied to a different quantity:
the number you report is a property of your protocol before it is a property of the air.

\item \textbf{Health-side records are all weak designs today, and one of them is
instructive about why.} Chin et al.\ regress monthly counts of HS clinical encounters on
monthly counts of days above 12\,$\mu$g/m$^3$ and report $\beta=0.57$ at lag 0 and
$\beta=0.62$ at lag 2, with lag 1 and lag 3 null. A non-monotonic lag structure across
four adjacent lags in a single-centre monthly series is a pattern of multiple comparisons,
not a biological latency.

\item \textbf{No forest plot again --- seventh consecutive issue.} Nothing in the corpus
is ratio-scale with an interval. The two quantitative health estimates are a
\textbf{$-$11.8\,g} birth-weight difference per 10\,$\mu$g/m$^3$ second-trimester
PM$_{2.5}$ and a \textbf{$+$23.1\%} preterm-birth risk increase, both quoted without CIs
in a narrative review. \texttt{f5\_forest.png} is intentionally absent.

\item \textbf{Standing gap, now measured over 469 records:} still not one paper couples an
instrument to a health endpoint. The MARTHA device flagged on 9 August
(\texttt{10.1016/j.jaerosci.2026.106882}) remains without an abstract in Crossref or
Europe PMC after 72 hours, so it cannot yet be counted.
\end{itemize}

\clearpage

% ==================================================================== ANALYTICS
\section*{\sffamily\bfseries\color{Deep}Corpus analytics}
\vspace{-1.5mm}{\color{Amber}\rule{\linewidth}{1.1pt}}\vspace{3mm}

\noindent\includegraphics[width=\linewidth]{\FIGDIR/f1_subtopics.png}
\figcap{Subtopic assignment of the 16 in-scope records. Assignment is single-label by
dominant contribution. Sensing/instrumentation (6) and exposure assessment (4) together
take 63\,\% of the issue; the six health-side subtopics carry one record each. The flat
tail is the shape of a day whose health literature was thin and whose measurement
literature was not.}

\vspace{3mm}
\noindent\includegraphics[width=\linewidth]{\FIGDIR/f2_design_endpoint.png}
\figcap{Study architecture (left) and endpoint class (right). Measurement campaigns lead
at \textbf{5}, with review/synthesis and modelling/inventory tied at \textbf{4} each and
one record each in experimental/toxicology, observational cross-sectional and
observational acute. On the right, \textbf{9 of 16 records carry no health endpoint at
all} --- 4 instrument, 4 exposure-surface, 1 algorithm --- and the remaining seven each
report a different endpoint, so no health outcome appears twice. That is the arithmetic
of a 63\,\% measurement corpus, and it is why this issue has no forest plot.}

\vspace{3mm}
\noindent\includegraphics[width=\linewidth]{\FIGDIR/f3_metric_geography.png}
\figcap{Particle metric (left) and geographic provenance (right). The left panel uses the
particle-metric vocabulary extended in this issue. It has not eliminated the fallback:
``PM (unspeciated) / emissions'' is still the largest single bin at \textbf{4}, but it now
holds only records that genuinely have no stated particle metric (an unspecified-pollutant
algorithm, unmeasured occupational dust, a qualitative pollution exposure and a
multi-pollutant-class review) rather than absorbing every instrument-side quantity as it
did through 9 August. \textbf{Optical properties} now carries \textbf{3} and
\textbf{composition / speciation} \textbf{2}, against \textbf{0} before the fix. On the
right, North America and the global/multi-region bucket tie at \textbf{4}; Europe has
\textbf{3}, China and East Asia \textbf{2} each, South Asia \textbf{1}. The
global/multi-region count includes \textbf{two} records that are honestly non-geographic
(a method review and an algorithm paper) plus two genuinely multi-region reviews, so it
should not be read as evidence of internationally pooled fieldwork.}

\vspace{3mm}
\noindent\includegraphics[width=\linewidth]{\FIGDIR/f6_lifecourse.png}
\figcap{Life-course coverage. One in-utero record (trimester-specific PM$_{2.5}$ and birth
weight), three adult, one older-adult, and \textbf{nothing at all in childhood or
adolescence} --- a gap that has now persisted across most issues of this watch and that
the 2--8 August weekly also recorded. Counts are of records, not participants.}

\vspace{3mm}
\noindent\includegraphics[width=\linewidth]{\FIGDIR/f4_heatmap.png}
\figcap{Subtopic $\times$ study architecture. The two blocks are disjoint, and that is the
finding. Sensing and exposure occupy measurement campaigns (3 and 2) and
modelling/inventory (2 and 2) and appear \textbf{nowhere} in the experimental or
observational columns; the six health-side subtopics occupy review/synthesis (3) and one
cell each in experimental/toxicology, observational cross-sectional and observational
acute, and appear \textbf{nowhere} under measurement campaign. No record in this issue
sits in both halves. Across all sixteen issues of this watch --- \textbf{469 records} ---
that intersection has never been populated, which makes it structural rather than a
sampling artefact of any one day.}

\clearpage

% ==================================================================== DIGEST
\section*{\sffamily\bfseries\color{Deep}Paper-level digest}
\vspace{-1.5mm}{\color{Amber}\rule{\linewidth}{1.1pt}}\vspace{2.5mm}

\footnotesize\noindent
Records are grouped by subtopic and ordered within group by first author. Tier A/B/C is a
judgement about how much weight the design can carry, not about the journal. Records
marked \textsc{metadata only} had no abstract in Crossref, Europe PMC, Semantic Scholar or
Scholar Gateway at build time and are listed rather than summarised. Records marked
\textsc{backfill} have entry dates before this watch began on 26 July 2026 and were
surfaced by a relevance-ranked connector, not by the date window.

\band{Sensing, forecasting \& instrumentation\hfill 6 records}{Amber}

\paperentry{Acton-Bond et al.\ 2026}
{Rolling positive matrix factorisation for aerosol optical apportionment}
{Atmos Environ}
{\textsc{metadata only} --- no abstract in Crossref, Europe PMC or Scholar Gateway.
Listed, not summarised. Carried because \emph{rolling} PMF applied to optical rather than
chemical data is a specific and consequential claim: conventional PMF source
apportionment requires filter sampling and offline speciation, whereas an optical
implementation with a rolling window could in principle run on an aethalometer or a
multi-wavelength absorption photometer at deployment cadence. That would move source
attribution from a laboratory capability to a network capability, which is the
distinction that decides whether a distributed sensing programme can report anything
beyond mass. Whether the paper supports that reading is unknown until an abstract
appears.}
{10.1016/j.atmosenv.2026.122295}

\paperentry{Akaezuwa et al.\ 2026}
{SFFO-RL intelligent environmental monitoring: smart sensing networks with artificial
intelligence}
{Environ Monit Assess}
{Proposes an ``Intelligent Environmental Monitoring Framework'' combining adaptive
filtering, Isolation Forest outlier removal, min--max normalisation, multi-order neighbour
feature fusion and a reinforcement-learning optimiser, reporting \textbf{97.1\,\%
accuracy, 97.8\,\% F1, 97.6\,\% recall and 98.21\,\% precision}. Those figures cannot be
interpreted. The abstract names no pollutant, no dataset, no sensor, no units, no sample
size and no held-out split; ``accuracy'' is quoted for what is elsewhere described as
predicting continuous pollutant concentrations, which is a category error unless the task
was silently discretised. A four-metric cluster in the 97--98\,\% band on environmental
time series with autocorrelation is the signature of evaluating on temporally adjacent
data. Carried as tier C because the framework class (edge-side denoising plus
learned correction) is squarely on this watch's axis and the record should be visible
in the register; the reported numbers should not be cited.}
{10.1007/s10661-026-15724-0}

\paperentry{Babu et al.\ 2026}
{Capturing the invisible: hyperlocal variability of size-resolved particulate matter in a
smart city}
{Environ Monit Assess}
{Mobile monitoring campaign in Madurai, Tamil Nadu, during the northeast-monsoon retreat:
a vehicle-mounted size-resolved optical instrument sampling every \textbf{6\,s} across
$\sim$\textbf{2.4\,km$^2$} aggregated to a \textbf{50\,m\,$\times$\,50\,m} grid.
Local enhancement above route background was largest for PM$_{10}$ at
\textbf{188.91\,$\pm$\,149.07\,$\mu$g/m$^3$}, attributed to road dust and construction;
evening PM$_{1}$ (\textbf{42.10\,$\mu$g/m$^3$}) exceeded morning levels; the size
distribution was \textbf{bimodal}, with coarse (16.77--32.37\,$\mu$m) and submicron
(0.275--0.53\,$\mu$m) peaks. Getis-Ord Gi* identified persistent hotspots on traffic
corridors and cold spots in restricted-access zones; SEM-EDX on 22 elements linked size to
composition. The limitation that threatens the inference is that the reported dispersion
(SD 79\,\% of the mean) is not decomposed: on a moving platform, true spatial variance,
uncalibrated optical response, inlet losses and 6\,s instrument response time are
confounded, and no co-location against a reference monitor is described. Directly relevant
to distributed sensing: this is the spatial resolution a mobile network buys, reported
without the calibration chain that would make it transferable.}
{10.1007/s10661-026-15746-8}

\paperentry{Bobbia et al.\ 2026}
{Multi-source dataset of NO$_2$ and PM from mobile and fixed sensors in Rouen (France)}
{Data in Brief}
{\textsc{backfill} (PubMed entry 2026-07-09; surfaced by the Amass connector, not by the
date window). Data descriptor for the TIGA ``Smart Mobility for All'' deployment in Rouen:
\textbf{10 mobile AtmoTrack} units on buses at 10\,s, \textbf{18 fixed AirSensEUR} units
on traffic lights at 1\,min, \textbf{6 reference stations} at 15\,min, plus 2 AtmoTrack
units co-located beside a reference station, over \textbf{15 January--30 November 2025}.
Mobile records carry bus speed, GPS position and onboard RH/T/pressure; fixed records add
O$_\mathrm{x}$, NO and CO. The stated purpose is explicitly \textbf{multi-hop
calibration} --- transferring a reference-derived correction through fixed sensors to
mobile units that never see a reference site --- and heterogeneous data fusion. This is
the single most directly usable record of the issue for distributed low-cost PM sensing,
because it is the rare case where the co-location topology needed to test calibration
transfer is published rather than described. Limitation: it is a data descriptor, so it
contains no calibration result, no error metric and no validation; the value is entirely
in the availability of the raw structure.}
{10.1016/j.dib.2026.113025}

\paperentry{Gilardoni et al.\ 2026}
{First continuous aerosol measurements at Testa Grigia at 3480\,m\,a.s.l.: aerosol
populations and dust transport dynamics in the southern European Alps}
{Atmos Chem Phys}
{First continuous \textbf{number size distribution and absorption coefficient} record at
1\,h resolution from the Testa Grigia Observatory (3480\,m\,a.s.l., Italian Alps),
spanning \textbf{September 2021--May 2023}. Three aerosol population types are separated
by transport pathway: a coarse population indicating long-range Saharan transport, a fine
population tied to Po Valley and alpine-valley boundary-layer dynamics and broader
continental flow, and low-number periods representing free-tropospheric background. The
record is a baseline dataset rather than a hypothesis test, and the limitation is the one
every single-site campaign carries --- population attribution rests on back-trajectory and
covariance reasoning, not on independent source measurement. Relevant to distributed
sensing for a narrow but useful reason: a high-altitude site with characterised free-
tropospheric periods gives an absolute-zero reference for optical instruments that
low-altitude urban co-location cannot provide, and the absorption-coefficient record is
one of the few instrument-side quantities published at deployment cadence.}
{10.5194/acp-26-11135-2026}

\paperentry{Wang et al.\ 2026}
{Measuring heavy metals in atmospheric particulate matter: operational definitions, method
comparability and environmental interpretation}
{Atmos Environ}
{\textsc{metadata only} --- no abstract in Crossref, Europe PMC or Scholar Gateway.
Listed, not summarised. Carried because the title states the problem this watch treats as
central, transposed to a different measurand: for heavy metals in PM, the reported
concentration depends on the digestion protocol, the filter substrate, the size cut and
whether total or bioaccessible fraction is meant --- so ``comparability'' is a property of
the method chain, not of the air. The same argument applies to low-cost optical mass
measurement, where the reported $\mu$g/m$^3$ is a function of assumed density, refractive
index and hygroscopic growth. A systematic treatment of operational definitions on the
metals side would be worth reading against the sensor-calibration literature.}
{10.1016/j.atmosenv.2026.122291}

\band{Exposure assessment \& modelling\hfill 4 records}{Clay}

\paperentry{Fan et al.\ 2026}
{Santa Ana winds control wildfire smoke transport and related health impacts in California
during January 2025}
{Atmos Environ}
{\textsc{metadata only} --- no abstract in Crossref, Europe PMC or Scholar Gateway.
Listed, not summarised. Carried because it is the first record in this watch to attach a
health-impact estimate to the January 2025 Los Angeles fires, and because the physical
premise is specific and testable: offshore Santa Ana flow reverses the usual coastal
transport direction, so smoke-attributable exposure is not predicted by proximity to the
burn perimeter. Satellite work on the same event (Li et al., \textit{AGU Advances}, Feb
2026) reports spread rates to \textbf{3.7\,km\,h$^{-1}$}, residential-area fire
intensities of \textbf{72--297\,W\,m$^{-2}$} and a total-PM emission coefficient of
\textbf{10.2--21.1\,g\,MJ$^{-1}$}, with 24\,h mean PM$_{2.5}$ within 7.5\,km reaching
\textbf{300\,$\mu$g/m$^3$} --- context, not this paper's result. The health-impact
methodology is unknown until an abstract appears, and for a wildland--urban-interface fire
burning structures the applicability of vegetation-fire concentration--response functions
is exactly the assumption that needs checking.}
{10.1016/j.atmosenv.2026.122292}

\paperentry{Gong et al.\ 2026}
{Regional analysis of urban physical factors influencing air pollution: evidence from
Seoul's living areas}
{Res Sq (preprint)}
{Constructs a \textbf{100\,m} grid over the Seoul metropolitan area, kriges a continuous
surface from \textbf{193 monitoring stations} for SO$_2$, CO, NO$_2$, O$_3$, PM$_{10}$ and
PM$_{2.5}$ combined into a Comprehensive Air-quality Index, reduces 15 urban-form
variables to four principal components (urban enclosure, low-density buildings,
high-density buildings, development intensity) and fits geographically weighted
regression. CAI rises from \textbf{88--92} in northeastern Seoul to \textbf{96--104} in
the southwest. The headline model comparison --- adjusted \textbf{R$^2$ 0.9928} for GWR
against \textbf{0.0182} for OLS --- should be read as a warning rather than a result. A
54-fold jump in explained variance on a kriged surface means the GWR is largely
reproducing the interpolation, since kriged values at 100\,m between 193 stations are
smooth by construction and locally varying coefficients will fit that smoothness almost
exactly. No spatial cross-validation or effective-degrees-of-freedom correction is
reported. Tier C. Carried because the urban-form-to-exposure framing is the right one for
siting a sensor network, and because it is a clean illustration of how interpolated
exposure surfaces manufacture model fit.}
{10.21203/rs.3.rs-10443285/v1}

\paperentry{Hernandez et al.\ 2026}
{Urbanization and humidity impact atmospheric deposition of anthropogenic microparticles
(including microplastics) in a humid subtropical city}
{Atmos Environ}
{\textsc{metadata only} --- no abstract in Crossref, Europe PMC or Scholar Gateway.
Listed, not summarised. Carried at tier C because atmospheric microplastic is a particle
class that optical monitors detect but do not identify: fibres in particular violate the
spherical-equivalent-diameter assumption underlying every optical particle counter, so a
deposition record that quantifies their abundance as a function of urbanisation and
humidity bears on the size and shape distribution a low-cost sensor is actually sampling.
The humidity dependence is the part worth reading if an abstract appears, since RH is
simultaneously the dominant artefact in optical PM measurement and, here, an exposure
variable.}
{10.1016/j.atmosenv.2026.122294}

\paperentry{Lin et al.\ 2026}
{Elevated aerosol layers within the daytime mixed layer over ecologically sensitive areas
of northwest China: diurnal variation, formation mechanisms and regional effects}
{Atmos Chem Phys}
{Identifies an elevated aerosol layer embedded \emph{within} the daytime mixed layer over
the Hexi Corridor using intensive vertical-profile campaigns in \textbf{2010 and 2012} at
Dunhuang and Minqin. Excluding dust events, elevated layers occurred at
\textbf{0.6--2\,km} on \textbf{82\,\%} of days at Dunhuang and \textbf{44\,\%} at Minqin,
with the paper analysing formation mechanism, air-quality consequence and radiative
effect. The relevance to surface monitoring is direct and under-appreciated: a
concentration maximum aloft but inside the mixing layer breaks the well-mixed assumption
that lets a surface monitor stand for a column, and it is precisely the configuration in
which satellite AOD and ground PM$_{2.5}$ decouple --- which is the failure mode of every
AOD-to-surface-PM scaling relationship. Limitations: two campaigns fourteen years old, two
sites, and an arid region whose boundary-layer behaviour does not generalise to the humid
urban settings where most low-cost networks are deployed.}
{10.5194/acp-26-11111-2026}

\band{Mechanistic toxicology\hfill 1 record}{Violet}

\paperentry{Cho et al.\ 2026}
{\textit{Lacticaseibacillus paracasei} ATG-E1 alleviates particulate-matter-induced airway
inflammation in association with altered inflammatory signalling, gut microbiota and
faecal metabolites}
{J Microbiol Biotechnol}
{BALB/c mice challenged intranasally with PM$_{10}$D (PM$_{10}$ plus diesel exhaust
particles) on days 0, 3, 6 and 8, then treated orally with ATG-E1 or dexamethasone
\emph{after} inflammation was established --- a therapeutic rather than prophylactic
design, which is the methodological point in its favour. ATG-E1 reduced BALF inflammatory
infiltrate and neutrophilia, lowered collagen deposition and pro-inflammatory mediators,
and attenuated \textbf{I$\kappa$B$\alpha$ and ERK phosphorylation} while leaving JNK and
p38 unchanged; caspase-1 and IL-1$\alpha$ fell, lung RNA-seq showed broad downregulation
of cytokine--cytokine-receptor signalling, and cecal 16S plus faecal metabolomics showed
microbiome remodelling with increased branched-chain fatty acids and decreased
branched-chain amino acids. Two limits: the microbiome and metabolite changes are
correlative, with no germ-free, antibiotic or transfer arm to establish that the gut axis
mediates the airway effect rather than co-occurring with it; and no effect size, n per
group or variance is given in the abstract, so the magnitude is unassessable. All authors
are employed by the company that owns the strain.}
{10.4014/jmb.2604.04029}

\band{Occupational \& indoor exposure\hfill 1 record}{Clay}

\paperentry{Zhao et al.\ 2026}
{Occupational dust exposure and thyroid nodule screening: a predictive modelling approach
for early detection in coal mining communities}
{J Int Med Res}
{Cross-sectional analysis of \textbf{1,708} coal miners examined at an occupational-disease
hospital in Anhui in April 2021, combining general, laboratory and imaging data to predict
nodular thyroid disease, using a proposed ``Nonlinear Inverse Nearest Manifold
Projection'' model benchmarked against established risk models. The record is carried at
tier C and should not be cited for a result: the abstract reports that the model
``achieved high scores in both F1 and area under the precision--recall curve'' and
``significantly outperformed'' comparators \textbf{without giving a single numeric value},
confidence interval, or comparator name. Beyond the missing numbers, dust exposure is
never measured --- occupation is the exposure proxy, with no mass concentration, no
respirable fraction, no duration and no job-exposure matrix --- so the paper cannot
support any dose--response claim about particulate exposure and thyroid nodules. Included
because occupational particulate exposure with thyroid endpoints is a thin literature and
the register should show what is in it.}
{10.1177/03000605261473853}

\band{Reproductive \& developmental\hfill 1 record}{Sky}

\paperentry{Hahad et al.\ 2026}
{Air pollution and adverse birth outcomes: a narrative review of epidemiological and
mechanistic findings}
{J Environ Health Sci Eng}
{\textsc{backfill} (PubMed entry 2026-06-29; surfaced by the Amass connector). Narrative
review synthesising epidemiological and mechanistic evidence linking prenatal air
pollution to low birth weight, preterm birth and stillbirth, with searches across PubMed,
Scopus and Web of Science. The two quoted estimates are a \textbf{$-$11.8\,g} birth-weight
difference and a \textbf{$+$23.1\,\%} increase in preterm-birth risk per
\textbf{10\,$\mu$g/m$^3$} second-trimester PM$_{2.5}$; mechanisms discussed are oxidative
stress, inflammation, endocrine disruption, vascular dysfunction and epigenetic
modification acting on placental morphology, perfusion and metabolic function. Both
estimates are quoted \textbf{without confidence intervals or provenance to a specific
pooled analysis}, which is why neither enters the forest plot --- a narrative review is
not a meta-analysis, and a point estimate lifted from one without its interval is not an
effect size. The review is honest about what remains unresolved: disentangling individual
pollutants, establishing causality, and identifying critical windows. Relevant to sensing
because trimester-specific windows demand exposure resolution that residential-address
annual means cannot deliver.}
{10.1007/s40201-026-00990-4}

\band{Neuro / mental health\hfill 1 record}{Deep}

\paperentry{Domingo 2026}
{Aging in a highly polluted world: challenges and solutions to prevent Alzheimer's disease}
{Arch Toxicol}
{\textsc{backfill} (PubMed entry 2026-06-09; surfaced by the Amass connector). Mechanistic
review drawing on human cohorts, animal models and \textit{in vitro} neuronal/glial
systems to link pollutant classes --- ambient PM, heavy metals, pesticides, PFAS and
micro/nanoplastics --- to Alzheimer's-relevant neuropathology, framed against the 2024
\textit{Lancet} Commission's inclusion of air pollution among 14 modifiable dementia risk
factors and a projected \textbf{$>$150\,million} cases by 2050. Shared pathways discussed
are neuroinflammation, oxidative stress, blood--brain-barrier disruption, tau-kinase
dysregulation, epigenetic reprogramming and gut--brain dysbiosis, with particular weight
on reports of AD-type neuropathology in children and young adults in heavily polluted
urban environments. That evidence is the review's most striking claim and also its
weakest link: it rests substantially on one research programme's Mexico City autopsy
series, where selection into autopsy, APOE distribution and co-exposure to metals and
traffic noise are not separable from PM. Read as an argument for exposure reduction as a
dementia-prevention lever, it is reasonable; read as evidence of a PM-specific causal
pathway, it is not yet that.}
{10.1007/s00204-026-04469-x}

\band{Burden, policy \& mitigation\hfill 1 record}{Coral}

\paperentry{Vartanyan et al.\ 2026}
{Structural inequities and chronic disease in Fresno County: a narrative review and
multi-sectoral framework}
{Cureus}
{Comparative descriptive analysis of public-health data for \textbf{2018--2021} showing
Fresno County, California, exceeding state prevalence and mortality for obesity,
hypertension, diabetes, asthma, coronary heart disease and stroke, attributed to poverty,
low educational attainment, provider shortages, food insecurity and severe air and
pesticide pollution, with disproportionate impact on racial and ethnic minority
populations. Tier C, and the reason is structural rather than a criticism of execution:
this is an ecological comparison of county aggregates in which air pollution is one of
several correlated exposures, none measured at individual level, so no exposure-specific
inference is available and none is claimed. Carried because the San Joaquin Valley is
among the most PM$_{2.5}$-burdened airsheds in the United States and because the paper's
underlying argument --- that monitoring density and exposure burden are inversely related
in marginalised communities --- is the same one Siu et al.\ quantified for eastern Canada
in the 5 August issue, here asserted rather than measured.}
{10.7759/cureus.112418}

\band{Other clinical endpoints\hfill 1 record}{Slate}

\paperentry{Chin et al.\ 2026}
{Impact of air quality and PM$_{2.5}$ on hidradenitis suppurativa activity in the
Sacramento region}
{Res Sq (preprint)}
{Retrospective analysis of \textbf{10,431} daily hidradenitis suppurativa encounters at UC
Davis Health, \textbf{April 2015--March 2022}, regressing monthly encounter counts on the
monthly count of days exceeding \textbf{12\,$\mu$g/m$^3$} PM$_{2.5}$ from the EPA Air
Quality System, with lags to three months. Positive associations at
\textbf{lag 0 ($\beta=0.5658$, $P=0.022$)} and \textbf{lag 2 ($\beta=0.6230$,
$P=0.015$)}; lag 1 and lag 3 null. The lag structure is the problem. A significant effect
at 0 and 2 months with nothing at 1 month has no biological reading, and four lags tested
without adjustment on a monthly series of $\sim$84 points is a multiple-comparisons
pattern. Three further limits: healthcare encounters measure care-seeking as much as
disease activity and are strongly seasonal; a regional monitor assigned to a referral
population is a substantial exposure misclassification; and no autocorrelation correction
or seasonal adjustment is described, which for monthly counts over seven years will
inflate significance on its own. Preprint, single centre, not peer reviewed.}
{10.21203/rs.3.rs-10570042/v1}

\clearpage

% ==================================================================== TABLE
\section*{\sffamily\bfseries\color{Deep}Corpus register}
\vspace{-1.5mm}{\color{Amber}\rule{\linewidth}{1.1pt}}\vspace{2.5mm}

{\fontsize{6.7}{8.4}\selectfont
\setlength{\tabcolsep}{2.4pt}
\begin{longtable}{@{}>{\RaggedRight}p{25mm}>{\RaggedRight}p{41mm}>{\RaggedRight}p{27mm}%
>{\RaggedRight}p{23mm}>{\RaggedRight}p{18mm}>{\RaggedRight}p{15mm}@{}}
\toprule
\textsf{\textbf{First author}} & \textsf{\textbf{Focus}} & \textsf{\textbf{Journal}} &
\textsf{\textbf{Design}} & \textsf{\textbf{Metric}} & \textsf{\textbf{Region}}\\
\midrule
\endfirsthead
\toprule
\textsf{\textbf{First author}} & \textsf{\textbf{Focus}} & \textsf{\textbf{Journal}} &
\textsf{\textbf{Design}} & \textsf{\textbf{Metric}} & \textsf{\textbf{Region}}\\
\midrule
\endhead
\multicolumn{6}{@{}l@{}}{\cellcolor{Mist}\textsf{\textbf{\textcolor{Deep}{Neuro / mental health}}}}\\[0.6mm]
\href{https://doi.org/10.1001/jamanetworkopen.2026.25015}{Dales} & Wildfire-sourced PM$_{2.5}$ and ED visits for migraine and other primary headache & \textit{JAMA Netw Open} & Case-crossover & PM$_{2.5}$ (wildfire) & Canada\\
\href{https://doi.org/10.1093/ageing/afag215}{Glover} & Umbrella meta-review of urban environmental exposures and cognitive health & \textit{Age Ageing} & Umbrella review / meta-analysis & PM$_{2.5}$ & Global\\
\href{https://doi.org/10.1016/j.nsa.2026.107019}{Goossens} & In utero air pollution exposure and pre-/postnatal brain development & \textit{Neurosci Appl} & Narrative review & PM$_{2.5}$ / PM$_{10}$ / PM$_{\mathrm{coarse}}$ & Global\\
\href{https://doi.org/10.1016/j.isci.2026.116617}{Ma} & PM$_{2.5}$-induced neurotoxicity and depression: exposure risks and mechanisms & \textit{iScience} & Narrative review & PM$_{2.5}$ + components & Global\\
\href{https://doi.org/10.3760/cma.j.cn112150-20250804-00752}{Pan} & Maternal PM exposure and offspring neurodevelopment: effects and mechanisms & \textit{Zhonghua Yu Fang Yi Xue Za Zhi} & Narrative review & PM (bulk) & China\\
\addlinespace[0.9mm]
\multicolumn{6}{@{}l@{}}{\cellcolor{Mist}\textsf{\textbf{\textcolor{Teal}{Cardiovascular \& metabolic}}}}\\[0.6mm]
\href{https://doi.org/10.3760/cma.j.cn112338-20251120-00843}{Huang} & PM$_{2.5}$ and dyslipidaemia in rural Chinese children; dietary effect modification & \textit{Zhonghua Liu Xing Bing Xue Za Zhi} & Cross-sectional (multistage) & PM$_{2.5}$ & China\\
\href{https://doi.org/10.32394/pe/218966}{Kazmierska} & Air pollution and cardiovascular disease: review of current data & \textit{Przegl Epidemiol} & Narrative review & PM$_{2.5}$ / PM$_{10}$ & Global\\
\href{https://doi.org/10.1007/s00125-026-06814-2}{Xi} & Air pollution, built environment and temperature vs. prediabetes phenotypes & \textit{Diabetologia} & Prospective cohort & PM$_{2.5}$ & Germany\\
\href{https://doi.org/10.1161/JAHA.125.046917}{Yang} & Short-term PM$_{2.5}$ and hospitalisation/ED use in atrial fibrillation & \textit{J Am Heart Assoc} & Case-crossover & PM$_{2.5}$ (wildfire-influenced) & USA\\
\addlinespace[0.9mm]
\multicolumn{6}{@{}l@{}}{\cellcolor{Mist}\textsf{\textbf{\textcolor{Sky}{Reproductive \& developmental}}}}\\[0.6mm]
\href{https://doi.org/10.1016/j.jes.2025.10.043}{Chen} & Prenatal air pollution, fetal growth trajectory phenotypes and child neurodevelopment & \textit{J Environ Sci (China)} & Prospective cohort & PM$_{2.5}$ / PM$_{10}$ & China\\
\href{https://doi.org/10.1016/j.ecoenv.2026.120566}{Huang} & Joint prenatal bisphenol + air-pollutant mixture exposure and preschool IQ & \textit{Ecotoxicol Environ Saf} & Prospective cohort (mixtures) & PM$_{2.5}$ / PM$_{10}$ & China\\
\href{https://doi.org/10.1007/s40572-026-00552-8}{Peterson} & Ambient air pollution, PAHs and pubertal development: critical review & \textit{Curr Environ Health Rep} & Critical review & PM$_{2.5}$ components & Global\\
\href{https://doi.org/10.1016/j.envres.2026.125313}{Xie} & Long-term PM$_{2.5}$ and timing of pubertal development in boys and girls & \textit{Environ Res} & Prospective cohort & PM$_{2.5}$ & China\\
\addlinespace[0.9mm]
\multicolumn{6}{@{}l@{}}{\cellcolor{Mist}\textsf{\textbf{\textcolor{Sage}{Respiratory \& allergic}}}}\\[0.6mm]
\href{https://doi.org/10.1007/s00484-026-03289-z}{Bayar Muluk} & Climate change and otolaryngology: airway disease impacts and emission reduction & \textit{Int J Biometeorol} & Narrative review & PM$_{2.5}$ & Global\\
\href{https://doi.org/10.4103/mgr.MEDGASRES-D-25-00164}{Liao} & Short-term air pollution, respiratory symptoms and lung function in asthma & \textit{Med Gas Res} & Panel study & PM$_{2.5}$ / PM$_{10}$ & China\\
\href{https://doi.org/10.3390/epidemiologia7040097}{Sanchez de Prada} & Ambient PM and asthma-coded admissions in children 0-3 y, Spain & \textit{Epidemiologia} & Ecological panel & PM$_{2.5}$ / PM$_{10}$ & Spain\\
\href{https://doi.org/10.1016/j.jhazmat.2026.143076}{Yang} & Allergenic vs. toxic risk of airborne seed hairs compared with ambient PM$_{2.5}$ & \textit{J Hazard Mater} & Field measurement & PM$_{2.5}$ + bioaerosol & China\\
\addlinespace[0.9mm]
\multicolumn{6}{@{}l@{}}{\cellcolor{Mist}\textsf{\textbf{\textcolor{Violet}{Mechanistic toxicology}}}}\\[0.6mm]
\href{https://doi.org/10.1016/j.jes.2025.12.074}{Ding} & PM$_{2.5}$ promotes vascular calcification via miR-27b-5p/HES1 signalling & \textit{J Environ Sci (China)} & In vitro + organoid + murine & PM$_{2.5}$ & China\\
\href{https://doi.org/10.1007/s00360-026-01698-5}{Maraschi} & Osmoregulatory disruption and metal bioaccumulation from settleable atmospheric PM & \textit{J Comp Physiol B} & Controlled exposure (ecotox) & Settleable APM & Brazil\\
\href{https://doi.org/10.1016/j.jes.2025.11.018}{Wu} & Macrophage NLRP3 inflammasome in real-ambient PM-induced hepatic glucose dysregulation & \textit{J Environ Sci (China)} & Real-ambient murine + co-culture & Real-ambient PM & China\\
\addlinespace[0.9mm]
\multicolumn{6}{@{}l@{}}{\cellcolor{Mist}\textsf{\textbf{\textcolor{Clay}{Occupational \& indoor}}}}\\[0.6mm]
\href{https://doi.org/10.3389/fmedt.2026.1796828}{Chen} & Nested suction instrument to cut intraoperative surgical-smoke PM$_{2.5}$ & \textit{Front Med Technol} & Bench / device experiment & PM$_{2.5}$ (surgical smoke) & China\\
\href{https://doi.org/10.1093/annweh/wxag055}{Cole} & Determinants of respirable crystalline silica exposure in Australian tunnelling & \textit{Ann Work Expo Health} & Observational exposure study & Respirable dust / RCS & Australia\\
\href{https://doi.org/10.1016/j.jhazmat.2026.143081}{Ding} & Stage-resolved hair Zn-isotope dynamics as a firefighting combustion-exposure tracer & \textit{J Hazard Mater} & Repeated-measures biomarker & Combustion-derived PM & China\\
\href{https://doi.org/10.3390/jox16040119}{Ramsden} & Xenobiotic hazards in aircraft cabin air & \textit{J Xenobiot} & Review + risk economics & Ultrafine particles & Global\\
\addlinespace[0.9mm]
\multicolumn{6}{@{}l@{}}{\cellcolor{Mist}\textsf{\textbf{\textcolor{Amber}{Sensing, forecasting \& instrumentation}}}}\\[0.6mm]
\href{https://doi.org/10.1016/j.jenvman.2026.130515}{Shi} & Adaptive spatiotemporal graph transformer for multi-site multi-step PM$_{2.5}$ forecasting & \textit{J Environ Manage} & Deep-learning model & PM$_{2.5}$ & China / India\\
\href{https://doi.org/10.3390/bios16070380}{Tsuchiya} & Seasonal and indoor-outdoor characteristics of fluorescent aerosol particles in offices & \textit{Biosensors} & Real-time sensor campaign & Bioaerosol / coarse & Japan\\
\addlinespace[0.9mm]
\multicolumn{6}{@{}l@{}}{\cellcolor{Mist}\textsf{\textbf{\textcolor{Amber}{Exposure assessment \& modelling}}}}\\[0.6mm]
\href{https://doi.org/10.1029/2025GH001621}{Solorzano-Araque} & Multiscale exposure modelling: LOTOS-EUROS + mobility surveys + particle toxicity & \textit{GeoHealth} & CTM + mobility integration & PM$_{2.5}$ / PM$_{10}$ & Colombia\\
\addlinespace[0.9mm]
\multicolumn{6}{@{}l@{}}{\cellcolor{Mist}\textsf{\textbf{\textcolor{Coral}{Burden, policy \& mitigation}}}}\\[0.6mm]
\href{https://doi.org/10.1002/hsr2.72886}{Karamzad} & Ischaemic heart disease burden in MENA 1990-2021; PM attributable fraction 28.2\% & \textit{Health Sci Rep} & GBD secondary analysis & Ambient PM & MENA\\
\href{https://doi.org/10.1371/journal.pone.0347989}{Rees} & HEPA filtration to reduce respiratory infection in care homes & \textit{PLoS One} & RCT process evaluation & Filtered PM & UK\\
\href{https://doi.org/10.1016/j.jes.2026.01.009}{Shan} & Policy-driven emission reductions and post-pandemic rebound, high-resolution inventories & \textit{J Environ Sci (China)} & Emission inventory & PM emissions & China\\
\addlinespace[0.9mm]
\multicolumn{6}{@{}l@{}}{\cellcolor{Mist}\textsf{\textbf{\textcolor{Slate}{Other clinical endpoints}}}}\\[0.6mm]
\href{https://doi.org/10.3390/dj14070440}{Lupita} & PM$_{2.5}$, preventive-care capacity and caries/edentulism burden across the EU & \textit{Dent J} & Ecological panel & PM$_{2.5}$ & EU-27\\
\href{https://doi.org/10.1016/j.jdin.2026.06.008}{Park} & Long-term PM exposure and risk of androgenetic alopecia & \textit{JAAD Int} & Nationwide cohort & PM$_{2.5}$ / PM$_{10}$ & South Korea\\
\href{https://doi.org/10.3171/2026.3.PEDS2615}{Sun} & Environmental exposure and area deprivation vs. paediatric hydrocephalus outcomes & \textit{J Neurosurg Pediatr} & Registry cohort + geospatial & PM$_{2.5}$ (satellite-derived) & USA\\
\addlinespace[0.9mm]
\bottomrule
\end{longtable}}


\vspace{4mm}
\begin{keybox}
{\sffamily\bfseries\footnotesize\color{Deep}Provenance and method}\\[1.2mm]
{\fontsize{7.2}{8.8}\selectfont
Entry window \textbf{2026-08-10 to 2026-08-10}, contiguous with the 9 August issue; no
gap. This is the first issue to run the \textbf{full connected connector set}, and the
per-leg yield is recorded here because the negative results are as informative as the
positive ones. \textbf{PubMed} connector, three independent queries against
\texttt{[Date - Entry]}: health axis 2, instrumentation axis 1, broad MeSH sweep 11 --- 13
distinct records, \textbf{6 carried}. Local harvester: PubMed health 5, PubMed sensing 0,
Europe PMC \texttt{CREATION\_DATE} 3 (\textbf{2 carried}), Crossref by ISSN across 17
journals 12 (\textbf{7 carried}). \textbf{Amass BiomedCore} returned 10 relevance-ranked
records: 4 were duplicates caught by \texttt{seen.json}, 3 were declined pre-watch
backfills, \textbf{3 carried as flagged backfill} with entry dates 9 July, 29 June and
9 June. \textbf{Consensus} returned 10, \textbf{0 carried} --- 4 were duplicates already
shipped on 26 July, 5 August and 8 August, the rest resolved outside the window; this is
the fourth consecutive run confirming that Consensus is a relevance index and cannot see a
single day. \textbf{Scholar Gateway} \texttt{semanticSearch} could not resolve any of the
five Elsevier abstracts: its corpus is Wiley full text with a stated last update of
\textbf{May 2026}, so it is a passage retriever for older literature, not a recency or
cross-publisher leg. \textbf{arXiv} 60 records, \textbf{1} published on or after 26 July
and off-topic. \textbf{medRxiv}, \textbf{Semantic Scholar} (rate-limited),
\textbf{Exa} and \textbf{Nimble} returned nothing inside the window. \textbf{OpenAlex}
HTTP 429, \textbf{eighth consecutive run}. \textbf{ClinicalTrials.gov} returned 311
matching trials but none registered or updated on 10 August; \texttt{trials.json}
unchanged, and the trial watch is surfaced in the weekly rollup, not here.
\textbf{27 records rejected}, each with a reason in \texttt{state/rejected.jsonl}.
\textbf{Five carried records have no abstract} in Crossref, Europe PMC, Semantic Scholar
or Scholar Gateway --- all five are Elsevier \textit{Atmospheric Environment} DOIs, the
same publisher-side gap first logged on 8 August, now persisting past 72 hours. They are
marked \textsc{metadata only}. Per the 9 August finding, the abstract re-check has been
moved to a weekly sweep.
\textbf{No forest plot is generated for this issue}: the corpus contains no ratio-scale
estimate with an interval, so \texttt{f5\_forest.png} is intentionally absent rather than
missing.
Two \texttt{plots.py} defects were closed in this issue. \texttt{pmclass()} gained four
instrument-side particle-metric classes (optical properties; size-resolved / number
distribution; composition / speciation; composite index) plus hyphen normalisation,
closing the defect logged on 9 August. Separately, \texttt{design\_group} fell through for
\textbf{12 of 16} records on the first build --- the ninth recurrence of that defect, and
the worst --- which rendered f2's architecture panel as a single ``Other / mixed'' block;
all twelve labels are now mapped and the figure was regenerated.
Subtopic, design, particle-metric and geography labels are assigned from title, abstract,
keywords and MeSH; assignment is single-label by dominant contribution and is a
judgement, not a controlled vocabulary.
All DOI links resolve to the publisher of record. Content derived from PubMed, Crossref,
Europe PMC and Amass metadata; abstracts remain the copyright of their respective
publishers.\par}
\end{keybox}

\end{document}
